% SHARD-ID: AISM-24-HCB3-DIAGONAL-INVERSE
% SHARD-TITLE: Diagonal Ha inverse propagation
% SHARD-SUMMARY: Reproduces lem-hcb3-diagonal-inverse, the af-validated propagation of level-one bijectivity and lower modulus to bijectivity of every amplification.
% SHARD-SUMMARY: Explains the exact entrywise inversion, the imported lower-modulus bound, and the elementary reciprocal estimate that produces the 1+C_inv*e inverse norm.
% SHARD-KEYWORDS: lem-hcb3-diagonal-inverse, af validated, Ha map, bijectivity, inverse norm, entrywise amplification
\section{Diagonal Ha inverse propagation}\label{sec:hcb3-diagonal-inverse}

\begin{lemma}[\texttt{lem-hcb3-diagonal-inverse}]\label{lem:hcb3-diagonal-inverse}
There are universal $C_{\mathrm{inv}}<\infty$ and $\comb_{\mathrm{inv}}>0$ such that, for
every H-CB datum with $\comb\le \comb_{\mathrm{inv}}$, \emph{if $\Ha{Q}{P}{P}$ has
level-one lower modulus at least $\tfrac14$ and is bijective at level one}, then every
amplification is bijective and
\begin{equation}\label{eq:di-main}
  \bigl\lVert\bigl(\Han{Q}{P}{P}{n}\bigr)^{-1}\bigr\rVert\;\le\;1+C_{\mathrm{inv}}\,\comb .
\end{equation}
\emph{Transcription note.} \textbf{Both} level-one hypotheses are load-bearing and neither
is discharged here: the conclusion is asserted only for those data whose level-one map is
\emph{already} bijective \emph{and} already has lower modulus at least $\tfrac14$. The two
are used at genuinely different points of the proof --- bijectivity supplies the inverse,
the lower modulus supplies its norm --- and neither implies the other in this setting. The
conditional shape is the corrected one, inherited from the same verdict that made the
lower-modulus clause of \texttt{conj-hcb} conditional after an unconditional inverse claim
was refuted by an exact $\mathbb C\oplus\mathbb C$ counterexample.
\end{lemma}

\noindent\textbf{Registry contract (verbatim).}
\contractquote{Diagonal Ha inverse propagation: there are universal C_inv < infinity and e_inv > 0 such that, for every H-CB datum with e <= e_inv, if Ha^Q_{P,P} has level-one lower modulus at least 1/4 and is bijective at level one, then every amplification is bijective and ||((Ha^Q_{P,P})_n)^(-1)|| <= 1+C_inv*e.}

\paragraph{Proof (prose account of the af-validated tree).}
The tree is short --- five nodes --- because the analytic work has already been done by
\texttt{lem-hcb3-diagonal-lower-modulus} (\S\ref{sec:hcb3-diagonal-lower-modulus}); what
remains is an exact algebraic step and one scalar inequality. Write $T:=\Ha{Q}{P}{P}$ and
$T_n:=\Han{Q}{P}{P}{n}$.

\emph{Constants.} Let $C_{\mathrm{diag}},\comb_{\mathrm{diag}}$ be universal witnesses
supplied by \texttt{lem-hcb3-diagonal-lower-modulus}. Enlarging $C_{\mathrm{diag}}$ only
weakens that external conclusion, so one may take $C_{\mathrm{diag}}>0$ and set
\[
  \comb_{\mathrm{inv}}:=\min\Bigl\{\comb_{\mathrm{diag}},\ \tfrac1{2C_{\mathrm{diag}}}\Bigr\},
  \qquad
  C_{\mathrm{inv}}:=2C_{\mathrm{diag}} .
\]
Both are universal, $\comb_{\mathrm{inv}}>0$, and the choice is made precisely so that
\begin{equation}\label{eq:di-half}
  0\le C_{\mathrm{diag}}\comb\le\tfrac12
  \qquad\text{whenever }0\le \comb\le \comb_{\mathrm{inv}} .
\end{equation}

\emph{Bijectivity amplifies exactly.} Fix an admissible datum and an arbitrary $n\ge1$.
This step uses the \emph{first} hypothesis and nothing else, and it involves no estimate
at all. By the meaning of matrix amplification, $T_n\bigl([Z_{ij}]\bigr)=[T(Z_{ij})]$ acts
slotwise. Surjectivity therefore passes upward: for any target $[Y_{ij}]$ the matrix
$[T^{-1}(Y_{ij})]$ is a preimage. So does injectivity: $T_n([Z_{ij}])=0$ forces
$T(Z_{ij})=0$ for every $i,j$, hence $Z_{ij}=0$ for every $i,j$ and $[Z_{ij}]=0$. Thus
$T_n$ is bijective and, moreover,
\begin{equation}\label{eq:di-entrywise}
  (T_n)^{-1}\ \text{is the entrywise amplification of }T^{-1} .
\end{equation}
Note what is \emph{not} claimed: nothing here bounds $\lVert(T_n)^{-1}\rVert$ ---
entrywise inversion alone gives no uniform norm control, which is why the second
hypothesis is needed.

\emph{The imported lower bound.} Since $\comb\le \comb_{\mathrm{inv}}\le \comb_{\mathrm{diag}}$
and the level-one lower modulus is at least $\tfrac14$ by the \emph{second} hypothesis,
\texttt{lem-hcb3-diagonal-lower-modulus} applies verbatim and gives
\begin{equation}\label{eq:di-lower}
  \nrm{T_n(Z)}\;\ge\;\bigl(1-C_{\mathrm{diag}}\comb\bigr)\nrm{Z}
  \qquad\text{for every }Z\in\Amp{n}{\Cnrs{P}} .
\end{equation}
This is the only place the $\tfrac14$ hypothesis enters, and it enters through the
imported lemma rather than being re-proved.

\emph{The reciprocal estimate.} Combine the two. For an arbitrary target element $Y$, put
$Z=(T_n)^{-1}Y$, which exists by \eqref{eq:di-entrywise}; then \eqref{eq:di-lower} reads
$\nrm{Y}=\nrm{T_n(Z)}\ge(1-C_{\mathrm{diag}}\comb)\nrm{Z}$. By \eqref{eq:di-half} the
quantity $x:=C_{\mathrm{diag}}\comb$ lies in $[0,\tfrac12]$, and on that interval the
elementary identity
\[
  \frac1{1-x}=1+\frac{x}{1-x}\;\le\;1+2x
\]
holds because $1-x\ge\tfrac12$. Therefore
\[
  \bigl\lVert(T_n)^{-1}Y\bigr\rVert
  \;\le\;\bigl(1-C_{\mathrm{diag}}\comb\bigr)^{-1}\nrm{Y}
  \;\le\;\bigl(1+2C_{\mathrm{diag}}\comb\bigr)\nrm{Y}
  \;=\;\bigl(1+C_{\mathrm{inv}}\comb\bigr)\nrm{Y}.
\]
Since the datum and $n$ were arbitrary, every amplification is bijective with the inverse
bound \eqref{eq:di-main}, and the constants are universal. The restriction to
$x\le\tfrac12$ is not cosmetic: it is what converts the reciprocal into a linear bound in
$\comb$, and it is exactly what the choice of $\comb_{\mathrm{inv}}$ secures.

\medskip
\noindent\textbf{Status.} \texttt{proved}; \texttt{af: validated} (T0). Workspace
\texttt{proofs/lem-hcb3-diagonal-inverse}: 5 nodes, all \texttt{validated}, root
\texttt{validated}, taint clean; first pass, two rounds, no challenges surviving.

\noindent\textbf{Role in the chain.} Imports
\texttt{lem-hcb3-diagonal-lower-modulus},
\texttt{lem-hcb2-amplified-adjointness},
\texttt{lem-hcb2-product-defect},
\texttt{lem-compcb-corner-algebra} and
\texttt{lem-hcb3-uniform-square-lower}. It completes the diagonal block of the H-CB tier:
with \S\ref{sec:hcb3-diagonal-upper-norm} and \S\ref{sec:hcb3-diagonal-lower-modulus} it
says that, \emph{under the two level-one hypotheses}, the amplified diagonal Ha map is a
bijection whose norm and inverse norm are both within $O(\comb)$ of $1$. The matching
off-diagonal statement \texttt{lem-hcb3-offdiagonal-inverse} is not validated, the hcb4
canonical tier is \texttt{af: none}, and the parent \texttt{conj-hcb} remains
\texttt{proved-mod-audit}. Nothing here bears on the status of \texttt{op-classical},
which is open.
